\begin{figure}[h]
\[
\begin{array}{llclr}
    \textrm{C types (CompCert subset)} 
     & \tau & ::=  & \mathrm{int}(n,\mathit{sg})           & \textrm{($n\!\in\!\{8,16,32,64\}$; signed/unsigned)} \\
     &      & \mid & \_Bool                                & \\
     &      & \mid & \mathrm{float}(p)                     & \textrm{($p\!\in\!\{32,64\}$)} \\
     &      & \mid & \tau [n]                              & \textrm{(array of $n$)} \\
     &      & \mid & \mathrm{struct}\ \{\, f_i : \tau_i \,\} & \textrm{(struct)} \\
     &      & \mid & \grey{ \mathrm{union}\ \{\, f_i : \tau_i \,\} } & \textrm{\grey{(unions not yet implemented)}} \\
     &      & \mid & \grey{\tau\,{*}}                             & \textrm{\grey{(pointers not yet implemented)}} \\
  \textrm{Field names} & & \ni & f & \\
\end{array}
\]
\caption{Cedar’s type universe, aligned with \textsc{CompCert} C.}
\label{fig:cedar-types}
\end{figure}

Figure~\ref{fig:cedar-typing} summarizes the core typing rules
that relate Cedar layouts ($\ell$) to C types ($\tau$).